\chapter{Quantum circuits}

\section{Quantum algorithms}
The promise of quantum computers is to enable new algorithms for the solution of problems that are intractable on classical computers (too many resources).
There are two broad classes of algorithm that are known to fulfill this promise.

The first class of algorithms is based on Shor's \emph{quantum Fourier transform},
and includes the factoring and discrete logarithm problem.
These algorithms provide an exponential speedup over the best known classical algorithms.

The second class of algorithms is based on Grover's algorithm for performing \emph{quantum searching}.
These provide a quadratic speedup over the best known classical algorithms for searching an unstructured database.

Coming up with good quantum algorithms appears to be a very difficult task.
Algorithm design is difficult,
and quantum algorithms need to be better than the best known classical algorithms to be useful.

\section{Single qubit operations}
A single qubit is a vector $\ket{\psi} = \alpha\ket{0} + \beta\ket{1}$, parametrized by two complex numbers $\alpha$ and $\beta$ satisfying $|\alpha|^2 + |\beta|^2 = 1$.
Operations on a qubit must preserve this norm, and thus are described by $2 \times 2$ unitary matrices.

The Pauli matrices are the following:
\[
    X = \begin{pmatrix}
        0 & 1 \\
        1 & 0
    \end{pmatrix},\quad
    Y = \begin{pmatrix}
        0 & -i \\
        i & 0
    \end{pmatrix},\quad
    Z = \begin{pmatrix}
        1 & 0 \\
        0 & -1
    \end{pmatrix}.
\]

Three other quantum gates are important for quantum computing: the Hadamard gate $H$, the phase gate $S$, and the $\pi/8$ gate $T$.
\[
    H = \frac{1}{\sqrt{2}}\begin{pmatrix}
        1 & 1 \\
        1 & -1
    \end{pmatrix},\quad
    S = \begin{pmatrix}
        1 & 0 \\
        0 & i
    \end{pmatrix},\quad
    T = \begin{pmatrix}
        1 & 0 \\
        0 & e^{i\pi/4}
    \end{pmatrix}.
\]

The Hadamard gate can be written in term of the Pauli matrices as
\[
    H = \frac{X + Z}{\sqrt{2}}.
\]

The Pauli matrices give raise to the \emph{rotation operators} about the $\hat{x}$, $\hat{y}$, and $\hat{z}$ axes when exponentiated:
\[
    R_x(\theta) 
    = e^{-i\theta X/2}
    = \cos\left(\frac{\theta}{2}\right)I - i\sin\left(\frac{\theta}{2}\right)X
    = \begin{pmatrix}
        \cos\frac{\theta}{2} & -i\sin\frac{\theta}{2} \\
        -i\sin\frac{\theta}{2} & \cos\frac{\theta}{2}
    \end{pmatrix},
\]
\[
    R_y(\theta)
    = e^{-i\theta Y/2}
    = \cos\left(\frac{\theta}{2}\right)I - i\sin\left(\frac{\theta}{2}\right)Y
    = \begin{pmatrix}
        \cos\frac{\theta}{2} & -\sin\frac{\theta}{2} \\
        \sin\frac{\theta}{2} & \cos\frac{\theta}{2}
    \end{pmatrix},
\]
\[
    R_z(\theta)
    = e^{-i\theta Z/2}
    = \cos\left(\frac{\theta}{2}\right)I - i\sin\left(\frac{\theta}{2}\right)Z
    = \begin{pmatrix}
        e^{-i\theta/2} & 0 \\
        0 & e^{i\theta/2}
    \end{pmatrix}.
\]
If $\hat{n} = (n_x, n_y, n_z)$ is a real unit vector in three dimensions,
a rotation by $\theta$ about the axis $\hat{n}$ is given by
\[
    R_{\hat{n}}(\theta) = e^{-i\theta \hat{n} \cdot \vec{\sigma}/2}
    = \cos\left(\frac{\theta}{2}\right)I - i\sin\left(\frac{\theta}{2}\right)\hat{n} \cdot \vec{\sigma}
\]
where $\vec{\sigma} = (X, Y, Z)$ is the vector of Pauli matrices.

An arbitrary unitary operator on a single qubit can be written in many ways as a combination of rotations,
together with a global phase shift on the qubit.

\begin{theorem}[$Z$-$Y$ decomposition for a single qubit]
    Suppose U is an unitary operation on a signle qubit.
    Then there exist real numbers $\alpha$, $\beta$, $\gamma$, and $\delta$ such that
    \[
        U = e^{i\alpha} R_z(\beta) R_y(\gamma) R_z(\delta).
    \]
\end{theorem}

\begin{corollary}\label{corollary:zy-decomposition}
    Suppose $U$ is a unitary gate on a single qubit.
    Then there exist unitary operators $A$, $B$, $C$ on a single qubit such that $ABC = I$ and $U = e^{i\alpha} A X B X C$ for some real number $\alpha$.
\end{corollary}

The following circuit identities are useful for simplifying quantum circuits:
\[
    HXH=Z, \quad HYH=-Y, \quad HZH=X.
\]

\section{Controlled operations}
The prototypical controlled operation is the controlled-NOT (CNOT) gate.
The two inputs of the CNOT gate are the \emph{control qubit} and the \emph{target qubit}.
In the computational basis, the action of the CNOT gate is given by $\ket{c, t} \to \ket{c, t \oplus c}$.
The matrix representation of the CNOT gate is
\[
    U_\text{CN} = \begin{pmatrix}
        1 & 0 & 0 & 0 \\
        0 & 1 & 0 & 0 \\
        0 & 0 & 0 & 1 \\
        0 & 0 & 1 & 0
    \end{pmatrix}.
\]
The circuit representation of the CNOT gate is the following:
\[
    \Qcircuit @C=1em @R=2em {
        & \ctrl{1} & \qw \\
        & \targ & \qw
    }
\]

A \emph{controlled-U} operation is a two-qubit operation (with control and target qubits)
that applies $U$ if the control qubit is set: $\ket{c}\otimes\ket{t} \to \ket{c}\otimes U\ket{t}$.
The circuit representation of a controlled-U operation is the following:
\[
    \Qcircuit @C=1em @R=2em {
        & \ctrl{1} & \qw \\
        & \gate{U} & \qw
    }
\]

The controlled-$U$ gate can be implemented using CNOT gates and single qubit operations.
Applying a phase shift to the target qubit is done with the following circuit:
\[
    \Qcircuit @C=1em @R=2em {
        & \ctrl{1} & \qw \\
        & \gate{
            \begin{pmatrix}
            e^{i\alpha} & 0 \\
            0 & e^{i\alpha}
            \end{pmatrix}
        } & \qw
    }
    =
    \Qcircuit @C=1em @R=2em {
        & \gate{
            \begin{pmatrix}
            1 & 0 \\
            0 & e^{i\alpha}
            \end{pmatrix}
        } & \qw \\
        & \qw & \qw
    }
\]
This corresponds to the $e^{i\alpha}$ factor in Corollary~\ref{corollary:zy-decomposition},
where $X$ is the NOT gate.
Using a CNOT gate instead of X,we have the following.
If the control qubit is set, $e^{i\alpha}AXBXC=U$ is applied to the target qubit.
If the control qubit is not set, $ABC=I$ is applied to the target qubit.
Therefore, the following circuit implements a controlled-$U$ operation:
\[
    \Qcircuit @C=1em @R=2em {
        & \ctrl{1} & \qw \\
        & \gate{U} & \qw
    }
    =
    \Qcircuit @C=1em @R=2em {
        & \qw & \ctrl{1} & \qw & \ctrl{1} & \gate{\begin{pmatrix} 1 & 0 \\ 0 & e^{i\alpha} \end{pmatrix}} & \qw \\
        & \gate{C} & \targ & \gate{B} & \targ & \gate{A} & \qw
    }
\]

More generally, a controlled-$U$ operation can be implemented with multiple control and target qubits.
Suppose we have $n+k$ qubits, and $U$ a $k$-qubit unitary operator.
We define the controlled operation $C^n(U)$ by the equation
\[
    C^n(U)\ket{x_1\cdots x_n}\ket{\psi} = \ket{x_1\cdots x_n}U^{x_1\cdots x_n}\ket{\psi}
\]
where $x_1\cdots x_n$ in the exponent of $U$ means the product of the bits $x_1,\ldots,x_n$.
The operator $U$ is applied to the last $k$ qubits if the fist $n$ qubits are all equal to one, otherwise nothing is done.
The circuit notation for $C^n(U)$ is the following:
\[
    \Qcircuit @C=1em @R=2em {
        & \ctrl{2} & \qw \\
        & \ctrl{1} & \qw \\
        & \multigate{1}{U} & \qw \\
        & \ghost{U} & \qw \\
    }
\]

Suppose $U$ is a single qubit unitary operator, and $V$ is an unitary operator so that $V^2 = U$.
Then the operation $C^2(U)$ can be implemented using the following circuit:
\[
    \Qcircuit @C=1em @R=2em {
        & \ctrl{2} & \qw \\
        & \ctrl{1} & \qw \\
        & \gate{U} & \qw \\
    }
    =
    \Qcircuit @C=1em @R=2em {
        & \qw & \ctrl{1} & \qw & \ctrl{1} & \ctrl{2} & \qw \\
        & \ctrl{1} & \targ & \ctrl{1} & \targ & \qw & \qw \\
        & \gate{V} & \qw & \gate{V^\dagger} & \qw & \gate{V} & \qw \\
    }
\]

The Toffoli gate is an especially important special case of $C^2(U)$, where $U$ is the NOT gate $X$: $C^2(X)$.
Defining $V = (1-i)(I + iX)/2$ and noting that $V^2 = X$,
we see that the circuit above implements the Toffoli gate using only one and two quibit operations.
One and two qubit reversible gates are sufficient to implement universal computation.
By contrast, one and two bits classical reversible gates are not sufficient to implement the Toffoli gate.

In fact, any unitary operation can be composed to an arbitrary good approximation from just the Hadamard, pase, controlled-NOT, and $\pi/8$ gates.
The Toffoli gate can be implemented using only these gates as follows:
\[
    \Qcircuit @C=1em @R=2em {
        & \ctrl{2} & \qw \\
        & \ctrl{1} & \qw \\
        & \targ & \qw \\
    }
    =
    \Qcircuit @C=.5em @R=2em {
        & \qw & \qw & \qw & \ctrl{2} & \qw & \qw & \qw & \ctrl{2} & \qw & \ctrl{1} & \qw & \ctrl{1} & \gate{T} & \qw \\
        & \qw & \ctrl{1} & \qw & \qw & \qw & \ctrl{1} & \qw & \qw & \gate{T^\dagger} & \targ & \gate{T^\dagger} & \targ & \gate{S} & \qw \\
        & \gate{H} & \targ & \gate{T^\dagger} & \targ & \gate{T} & \targ & \gate{T^\dagger} & \targ & \gate{T} & \gate{H} & \qw & \qw & \qw & \qw \\
    }
\]

$C^n(U)$ can be implemented using $n-1$ working qubits (in addition to the $n$ control qubits and the target qubit),
which all start and end in the $\ket{0}$ state.
Suppose that control qubits are in the computational basis state $\ket{c_1,\ldots, c_n}$.
The first stage of the circuit is to reversibily AND all the control qubits together to produce $c_1\cdot \ldots \cdot c_n$ in the first working qubit,
by applying Toffoly gates, until the final working qubit is in the state $\ket{c_1\cdot \ldots \cdot c_n}$.
Next, the $U$ operation on the target qubit is applied,
conditional on the final working qubit being set to one.
Finally, the last part of the circuit reverses the steps of the first stage,
to return all the working qubits to the $\ket{0}$ state.
The combined result is to apply $U$ to the target qubit,
if and only if all the control bits $c_1,\ldots, c_n$ are set to one.

For $n = 5$, the circuit is as follows:
\[
    \Qcircuit @C=1em @R=2em {
        \lstick{\ket{c_1}} & \ctrl{5} & \qw & \qw & \qw & \qw & \qw & \qw & \qw & \ctrl{5} & \qw \\
        \lstick{\ket{c_2}} & \ctrl{4} & \qw & \qw & \qw & \qw & \qw & \qw & \qw & \ctrl{4} & \qw \\
        \lstick{\ket{c_3}} & \qw & \ctrl{4} & \qw & \qw & \qw & \qw & \qw & \ctrl{4} & \qw & \qw \\
        \lstick{\ket{c_4}} & \qw & \qw & \ctrl{4} & \qw & \qw & \qw & \ctrl{4} & \qw & \qw & \qw \\
        \lstick{\ket{c_5}} & \qw & \qw & \qw & \ctrl{4} & \qw & \ctrl{4} & \qw & \qw & \qw & \qw \\
        \lstick{\ket{0}} & \targ & \ctrl{1} & \qw & \qw & \qw & \qw & \qw & \ctrl{1} & \targ & \qw \\
        \lstick{\ket{0}} & \qw & \targ & \ctrl{1} & \qw & \qw & \qw & \ctrl{1} & \targ & \qw & \qw \\
        \lstick{\ket{0}} & \qw & \qw & \targ & \ctrl{1} & \qw & \ctrl{1} & \targ & \qw & \qw & \qw \\
        \lstick{\ket{0}} & \qw & \qw & \qw & \targ & \ctrl{1} & \targ & \qw & \qw & \qw & \qw \\
        \lstick{\ket{t}} & \qw & \qw & \qw & \qw & \gate{U} & \qw & \qw & \qw & \qw & \qw \\
    }
\]

An open circle notation indicate conditioning on the qubit set to zero, while a closed circle indicates conditioning on the qubit being set to one.
\[
    \Qcircuit @C=1em @R=2em {
        & \ctrlo{1} & \qw \\
        & \targ & \qw \\
    }
    =
    \Qcircuit @C=1em @R=2em {
        & \gate{X} & \ctrl{1} & \gate{X} & \qw \\
        & \qw & \targ & \qw & \qw \\
    }
\]

\section{Measurement}
Measurement in quantum circuits is denoted by a meter symbol,
which represents a projective measurement in the computational basis.
It is conventional to not use any special symbol for more general measurements,
they  can always be represented by unitary transformations with ancilla qubits,
followed by projective measurements.

There are two important principles about measurement in quantum circuits:
\begin{description}
    \item[Principle of deferred measurement] Measurements can always be moved from an intermediate stage of a quantum circuit to the end of the circuit;
        if the measruement results are used at any stage of the circuit,
        then the classically controlled operations can be replaced by conditional quantum operations.
    \item[Principle of implicit measurement] Without loss of generality,
        any unterminated quantum wires (qubits with are not measured) at the end of a quantum circuit may be assumed to be measured.
\end{description}

Quantum measurements are often performed as an intermediate step in a quantum circuit,
and the measurement results are used to conditionally control subsequent quantum gates.
This was the case, for example, in the quantum teleportation circuit.
The principle of deferred measurement, applied to quantum teleportation, results in the following circuit:
\[
    \Qcircuit @C=1em @R=2em {
        \lstick{\ket{\psi}} & \qw & \ctrl{1} & \gate{H} & \qw & \ctrl{2} & \meter  \\
        \lstick{} & \qw & \targ & \qw & \ctrl{1} & \qw & \meter \\
        \lstick{} & \qw & \qw & \qw & \gate{X} & \gate{Z} & \qw & \rstick{\ket{\psi}} \qw 
        \inputgroupv{2}{3}{0.6em}{1.75em}{\ket{\beta_{00}}}  \\ 
    }
\]

A consequence of the principle of deferred measurement is that \emph{measurements commute with controls}:
\[
    \Qcircuit @C=1em @R=2em {
        & \qw & \ctrl{1} & \meter & \qw \\
        & \qw & \gate{U} & \qw & \qw \\
    }
    =
    \Qcircuit @C=1em @R=2em {
        & \qw & \meter & \cw & \\
        & \qw & \qw & \gate{U} \cwx[-1] & \qw \\
    }
    \equiv
    \Qcircuit @C=1em @R=2em {
        & \qw & \meter & \\
        & \qw & \gate{U} \cwx[-1] & \qw \\
    }
\]

Consider a circuit containing just two qubit, with the first qubit measured at the end of the circuit.
The principle of implicit measurement originates from the fact that the reduced density matrix of the first qubit is not affected by performing a measurement on the second qubit.

\section{Universal quantum gates}
In classical computation,
a small set of gates (e.g. AND, OR, NOT),
said to be \emph{universal},
can be used to compute any arbitrary classical function.

In quantum computation,
a set of gates is said to be \emph{universal for quantum computation} if any unitary operation may be approximated to arbitrary accuracy by a quantum circuit involving only those gates.

\subsection{Two-level unitary gates are universal}
Arbitrary unitary operator may be expressed \emph{exactly} as a product of unitary operators tat each acts non-trivially only on a subspace spanned by two computational basis states.
A unitary matrix $U$ can be decomposed into a product of \emph{two-level unitary matrices}.

Suppose $U$ acts on a $d$-dimensional space.
Then,
we can find two-level unitary matrices $U_1, \ldots, U_{d-1}$ such that the matrix $U_{d-1}\cdots U_1 U$ has all zero in the first row and column except for the first element,
which is equal to one.
We then repeat this procedure for the $(d-1) \times (d-1)$ unitary submatrix in the lowe right hand corner,
and so on,
with the end result that an arbitrary $d \times d$ unitary matrix can be decomposed as
\[
    U = V_1 V_2 \cdots V_k
\]
where the matrices $V_i$ are two-level unitary matrices,
and $k \le (d-1) + (d-1) + \cdots + 1 = d(d-1)/2$.

A corollary of the above is that an arbitrary unitary matrix on an $n$ qubit system may be written as a product of at most $2^{n-1}(2^n-1)$ two-level unitary matrices.

\subsection{Single qubit and CNOT gates are universal}
An arbitrary unitary operator may be expressed \emph{exactly} using single qubit and CNOT gates.
This follows from the fact that any two-level unitary matrix can be implemented using single qubit and CNOT gates.

Suppose we have two distinct binary numbers $s$ and $t$ of length $n$.
A \emph{Gray code} connecting $s$ and $t$ is a sequence of binary numbers,
starting with $s$ and ending with $t$,
such that adjacent members of the list differ in exactly one bit.\graffito{
A Gray code connecting $s=101001$ and $t=110011$ is
\[
    \begin{matrix}
        1 & 0 & 1 & 0 & 0 & 1 \\
        1 & 0 & 1 & 0 & 1 & 1 \\
        1 & 0 & 0 & 0 & 1 & 1 \\
        1 & 1 & 0 & 0 & 1 & 1 \\
    \end{matrix}
\] 
}
We can always find a Gray code such that $m \le n + 1$ since $s$ and $t$ differ in at most $n$ locations.

Let $g_1, \ldots, g_m$ be the elements of a Gray code connecting $s$ and $t$,
with $g_1 = s$ and $g_m = t$.
The idea of a quantum circuit implementing U is to perform a sequence of gates effecting the state changes $\ket{g_1} \to \ket{g_2} \to \cdots \to \ket{g_m}$,
then to perform a controlled-$\tilde{U}$ operation with the target qubit located at the single bit where $g_{m-1}$ and $g_m$ differ,
and then to undo the first sequence of gates,
transforming $\ket{g_{m-1}} \to \ket{g_{m-2}} \to \cdots \to \ket{g_1}$.
 
The implementation of a two-level unitary operation $U$ requires at most $2(n-1)$ controlled operations to swap $\ket{g_1}$ with $\ket{g_{m-1}}$ and back again.
Each of these controlled operations can be implemented using $O(n)$ single qubit and CNOT gates.
The controlled $\tilde{U}$ gate also requires $O(n)$ gates.
Thus, implementing $U$ requires $O(n^2)$ single qubit and CNOT gates.
An arbitrary unitary operation on $n$ qubits can be implemented using $O(n^24^n)$ single qubit and CNOT gates.

\subsection{A discrete set of universal operations}
A discrete set of gates can't be used to implement an arbitrary unitary operation exactly,
since the set of unitary operations is continuous.
However, a discrete set can be used to \emph{approximate} any unitary operation.

Let $U$ be the target unitary operation to be implemented, and let $V$ be the unitary operation that is actually implemented.
The \emph{error} of the implementation is defined as
\[
    E(U, V) = \max_{\ket{\psi}} \norm{(U - V)\ket{\psi}},
\]
where the $\max$ is taken over all normalized quantum states $\ket{\psi}$ in the state space.
If $E(U, V)$ is small, any measurement performed on the state $V\ket{\psi}$ will yield approximately the same measurement statistics as measurements of $U\ket{\psi}$.

If $M$ is a POVM element in an arbitrary POVM,
and $P_U$ and $P_V$ are the probabilities of obtaining the measurement outcome corresponding to $M$ when measuring $U\ket{\psi}$ and $V\ket{\psi}$ respectively,
then
\[
    \abs{P_U - P_V} \leq 2E(U, V).
\]
Therefore, if $E(U, V)$ is small,
then measurement outcomes occur with similar probabilities regardless of whether $U$ or $V$ is performed.

If we perform a sequence of gates $V_1, \ldots, V_m$ to implement some other sequence of gates $U_1, \ldots, U_m$,
then the error of the implementation is bounded by
\[
    E(U_m\cdots U_1, V_m\cdots V_1) \leq \sum_{j=1}^m E(U_j, V_j).
\]
The errors add at most linearly with the number of gates in the sequence.

We can consider different set of gates that are universal for quantum computation:
the \emph{standard set} of universal gates consists of the Hadamard gate, the thphase gate, the controlled-NOT gate, and the $\pi/8$ gate,
while another set of universal gates consists of the Hadamard gate, the phase gate, the controlled-NOT gate, and the Toffoli gate.

The Hadamard and $\pi/8$ gates can be used to approximate any single qubit unitary operation to arbitrary accuracy.
This is because $T$ (a rotation by $\pi/4$ radians about the $\hat{z}$ axis on the Block sphere)
and $HTH$ (a rotation by $\pi/4$ radians about the $\hat{x}$ axis on the Bloch sphere)
can be used to construct a rotation $R_{\hat{n}}(\theta)$,
and a repeated application of $R_{\hat{n}}(\theta)$ can be used to approximate a rotation by any angle.
Since
\[
    E\left(U, R_{\hat{n}}^{n_1}(\theta)HR_{\hat{n}}^{n_2}(\theta)HR_{\hat{n}}^{n_3}(\theta)\right) < \epsilon,
\]
given any single qubit unitary operator $U$ and any $\epsilon > 0$,
it is possible to approximate $U$ to accuracy $\epsilon$ using a circuit composed of Hadamard and $\pi/8$ gates.

Given a quantum circuit containing $m$ gates, either CNOTs or single qubit unitary gates,
we may approximate it using Hadamard, controlled-NOT, and $\pi/8$ gates
(the phase gate allows to make the approximation fault-tolerant).
To achieve an overall error of $\epsilon$,
we need to approximate each gate in the circuit to an accuracy of $\epsilon/m$.
Approximating the $m$ gate circuit requires $\Omega(m2^{m/\epsilon})$ gates.
A much better rate of convergence can be achieved using the \emph{Solovay-Kitaev theorem},
which implies that an arbitrary single qubit gate may be approximated to an accuracy $\epsilon$ using $O(\log^c(1/\epsilon))$ gates from a discrete universal set of gates,
where $c \approx 2$.

The Hadamard, phase, controlled-NOT, and $\pi/8$ gates are universal for (fault-tolerant) quantum computation;
given a circuit containing CNOTs and single qubit unitary gates,
it is possible to simulate the circuit to good accuracy,
and the simulation can be performed efficiently (polynomial in $\log(m/\epsilon)$).

\subsection{Approximating arbitrary unitary gates is generically hard}
Given a unitary transformation $U$ on $n$ qubits, there is not always a circuit of size polynomial in $n$ approximating $U$.
There are unitary operations requiring exponentially many operations.

\section{Summary of the quantum circuit model of computation}
The following are the key elements of the quantum circuit model of computation:
\begin{description}
    \item[Classical resources] A quantum computer consists on a classical part and a quantum part.
        In principle, there is no need for the classical part of the computer,
        but in practice certain tasks may be made much easier if parts of the computation can be done classically.
    \item[A syuitable state space] A quantum circuit operates on some number $n$ of qubits.
        The state space is thus a $2^n$-dimensional Hilbert space.
        Product states of the form $\ket{c_1, \ldots, c_n}$ are called \emph{computational basis states} of the computer.
    \item[Ability to prepare statees in the computational basis] It is assumed that any computational basis state $\ket{x_1, \ldots, x_n}$ can be prepared in at most $n$ steps.
    \item[Ability to perform quantum gates] Gates can be applied to any subset of qubits as desired,
        and a universal family of gates can be implemented.
        The Hadamard, phase, controlled-NOT, and $\pi/8$ gates form a family of gates from which any unitary operation can be approximated.
    \item[Ability to perform measurements in the computational basis] Measurements may be performed in the computational basis of one or more of the qubits in the computer.
\end{description}

The ``quantum Turing machine'',
a quantum generalization of the classical Turing machine model,
has been shown to be equivalent to the quantum circuit model of computation.

\section{Simulation of quantum systems}
One of the most important practical applications of computation is the simulation of physical systems.

\subsection{Simulation in action}
Simulation of quantum systems by classical computers is possible, but generally inefficient.
The dynamical behavior of a quantum system is described by the time-dependent Schrödinger equation,
\[
    i\hbar \frac{d}{dt}\ket{\psi(t)} = H\ket{\psi(t)}.
\]
Using the position representation $\psi(x, y) = \braket{x| \psi(t)}$, the Schrödinger equation can be written as
\[
    i\hbar \frac{\partial}{\partial t}\psi(x, t) = \left(-\frac{1}{2m}\frac{\partial^2}{\partial x^2} + V(x)\right)\psi(x, t).
\]

The key challenge in quantum systems is the exponential number of differential equations which must be solved.
For $n$ qubits evolving according to the Schrödinger equation, there are $2^n$ coupled differential equations to solve.

Quantum computers can efficiently simulate quantum systems for which there is no known efficient classical algorithm.

\subsection{The quantum simulation algorithm}
The formal solution to the time-dependent Schrödinger equation is given by
\[
    \ket{\psi(t)} = e^{-iHt/\hbar}\ket{\psi(0)}.
\]
In principle, $H$ is extremely difficult to exponentiate,
and a first-order approximation $e^{-iHt/\hbar} \approx I - iHt/\hbar$ is generally not satisfactory.

In practice, many Hamiltonians (Hubbard model, Ising model, \dots) can be written as a sum of (local) terms:
\[
    H = \sum_{k=1}^L H_k.
\]
While $e^{-iHt/\hbar}$ is difficult to compute, $e^{-iH_k t/\hbar}$ acts on a much smaller subsystem,
and is straightforward to approximate using a quantum circuit.
However, in general, $[H, H_k] \neq 0$, and thus $e^{-iHt/\hbar} \neq \prod_{k=1}^L e^{-iH_k t/\hbar}$.

\begin{theorem}[Trotter formula]
    Let $A$ and $B$ be Hermitian operators.
    Then, for any $t$,
    \[
        \lim_{n \to \infty} \left(e^{-iAt/n}e^{-iBt/n}\right)^n = e^{-i(A+B)t}.
    \]
\end{theorem}
The Trotter formula is true even if $A$ and $B$ do not commute.

The \emph{Baker-Campbell-Hausdorff formula} tells us that
\[
    e^{(A+B)\Delta t} = e^{A\Delta t}e^{B\Delta t}e^{-\frac{1}{2}[A, B]\Delta t^2} + O(\Delta t^2).
\]

\begin{algorithm}
    \caption{Quantum simulation}
    \begin{algorithmic}[1]
        \Require
        A Hamiltonian $H = \sum_{k=1}^L H_k$ acting on an $N$-dimensional system,
        where each $H_k$ acts on a small subsystem of the size independent of $N$.
        
        \Require 
        An initial state $\ket{\psi_0}$ of the system at time $t=0$.
        
        \Require
        A positive non-zero accuracy $\delta$.
        
        \Require
        A time $t_f$ at which te evolved system is desired.

        \Ensure
        A state $\ket{\tilde{\psi}(t_f)}$ such that $\abs{\ket{\tilde{\psi}(t_f)} - e^{-iHt_f}\ket{\psi_0}}^2 \ge 1 - \delta$.

        \Statex
        \textbf{Runtime:} $O(\text{poly}(1/\delta))$ operations.

        \Statex
        \textbf{Procedure:} Choose a representation such that the state $\ket{\tilde{\psi}}$ of $n=\text{poly}(\log N)$ qubits approximates the system and the operators $e^{-iH_k t}$ can be implemented efficiently on a quantum computer.
        Select an approximation method and $\Delta t$ such that the expected error is acceptable (and $j\Delta t =t_f$ for some integer $j$),
        construct the corresponding quantum circuit $U_{\Delta_t}$ for the iterative step and do:

        \State 
        $\ket{\tilde{\psi}_0} \gets \ket{\psi_0}$; $j=0$ \Comment{Initialize state}

        \State 
        $\to \ket{\tilde{\psi}_{j+1}} \gets U_{\Delta t}\ket{\tilde{\psi}_j}$ \Comment{Iterative update}

        \State
        $\to j = j + 1$; goto 2 until $j\Delta t\ge t_f$ \Comment{Loop}

        \State
        $ \to \ket{\tilde{\psi}(t_f)} = \ket{\tilde{\psi}_j}$ \Comment{Final result}
    \end{algorithmic}
\end{algorithm}

\subsection{An illustrative example}
Quantum simulations are possible even for Hamiltonians which act non-trivially on all or nearly all parts of a large system.
Suppose
\[
    H = Z_1 \otimes Z_2 \otimes \cdots \otimes Z_n,
\]
which acts on an $n$-qubit system.
What we need is a quantum circuit which implements $e^{-iH\Delta t/\hbar}$.

Although the Hamiltonian $H$ involves all the qubits in the system,
it does so in a classical manner:
the phase shift applied to the system is $e^{-i\Delta t/\hbar}$ if the \emph{parity} of the $n$ qubits is even;
otherwise the phase shift is $e^{i\Delta t/\hbar}$.
It is possible to simulate $H$  by first computing the parity and storing the result in an ancilla qubit,
then applying the appropriate phase shift conditioned on the parity,
then uncomputing the parity (erase the ancilla qubit).

For $n=3$ we can use the following circuit:
\[
    \Qcircuit @C=1em @R=2em {
        \lstick{} & \ctrl{3} & \qw & \qw & \qw & \qw & \qw & \ctrl{3} & \qw \\
        \lstick{} & \qw & \ctrl{2} & \qw & \qw & \qw & \ctrl{2} & \qw & \qw \\
        \lstick{} & \qw & \qw & \ctrl{1} & \qw & \ctrl{1} & \qw & \qw & \qw \\
       \lstick{\ket{0}} & \targ & \targ & \targ & \gate{e^{-iZ\Delta t/\hbar}} & \targ & \targ & \targ & \rstick{\ket{0}}\qw \\
    }
\]

Extending the same procedure, we can efficiently simulate any Hamiltonian of the form
\[
    H = \bigotimes_{k=1}^n \sigma^k_{c(k)},
\]
where $\sigma^k_{c(k)}$ is a Pauli matrix (or the identity matrix) acting on the $k$-th qubit,
and $c(k) \in \{0, 1, 2, 3\}$ specifies which Pauli matrix is used ($I$, $X$, $Y$, or $Z$).
The qubits upon which the identity operates can be ignored,
and the $X$ and $Y$ operations can be transformed by single qubit gates to $Z$ operations.

It is possible to simulate any Hamiltonian of the form
\[
    H = \sum_{k=1}^L H_k,
\]
where the only restriction is that the individual terms $H_k$ have a tensor product structure,
and than L is polynomial in the number of qubits $n$.
More generally, all that is required is that the individual terms $H_k$ can be efficiently simulated on a quantum computer.
